% ============================================================
% CHAPTER 11
% ============================================================

\chapter{Derived Mapping Stacks and Scattering}
\label{chap:derived-scattering}

\section{Motivation}

Derived mapping stacks provide a powerful framework for describing
field configurations as morphisms between derived geometric objects. In
the lamphrodynamic setting, scattering processes may be understood as
derived morphisms between boundary configurations, with amplitudes
encoded by mapping spaces in the derived category.

This chapter outlines a speculative AKSZ--style construction for
lamphrodynamic scattering, emphasizing conceptual connections rather
than detailed computations.


\section{Derived Mapping Stacks}

Given derived stacks $X$ and $Y$, the mapping stack $\Map(X,Y)$
parameterizes morphisms $X\to Y$ up to homotopy. Scattering processes
correspond to derived maps from incoming to outgoing boundary
configurations.

\begin{definition}
The \emph{lamphrodynamic scattering stack} is
\[
  \Sc := \Map(B_{\mathrm{in}},B_{\mathrm{out}})
\]
where $B_{\mathrm{in}},B_{\mathrm{out}}$ are derived lamphrodynamic
boundary stacks.
\end{definition}

\begin{remark}
$\Sc$ generalizes traditional scattering amplitudes by incorporating
derived geometry and lamphrodynamic structure.
\end{remark}


\section{AKSZ--like Construction}

Analogous to Volume~I, Chapter~7, an AKSZ--type construction yields a BV
action
\[
  S_{\mathrm{sc}} \colon \Map(B_{\mathrm{in}},B_{\mathrm{out}})\to\mathbb{R}
\]
whose critical points describe scattering processes.

\begin{proposition}[Heuristic]
Assuming lamphrodynamic boundary stacks carry a $(-1)$--shifted
symplectic structure, the BV action $S_{\mathrm{sc}}$ satisfies the
classical master equation $\{S_{\mathrm{sc}},S_{\mathrm{sc}}\}=0$.
\end{proposition}


\section{Speculative Directions}

\begin{enumerate}
\item Entropy--dependent deformation of scattering geometry.
\item Relation to ambitwistor scattering integrals.
\item Potential amplituhedron--like lamphrodynamic polytopes.
\end{enumerate}


\section{Conclusion}

While speculative, these ideas connect lamphrodynamic field theory to
modern scattering formalisms, suggesting deep geometric relationships
between entropy, derived stacks, and holomorphic scattering structures.

\medskip
\noindent
This completes Part~V. We now turn to creative and philosophical
dimensions.

\newpage
